% =====================================================================
% 応用編 第1章 Issueで作業を管理する
% ステータス: 雛形（見出しと執筆メモのみ。本文はこれから書く）
% =====================================================================
\chapter{Issueで作業を管理する}

別冊『チームで安全に開発する』では、\textbf{ブランチを分けてPRでマージする}という「変更の流し方」を扱いました。この章で扱うのは、その1つ手前にある\textbf{「何を、誰が、いつやるのか」}の管理です。GitHubでは、これを\textbf{Issue（イシュー）}という機能で行います。

\begin{draftbox}
\textbf{この章のゴール}（清書時にこの枠は削除する）
\begin{itemize}[leftmargin=1.4em, nosep]
  \item Issueを立てて、作業をチームで見える化できる
  \item ラベル・担当者・マイルストーンで、Issueを整理できる
  \item IssueとPRを結びつけ、マージと同時にIssueを閉じられる
  \item Projects（ボード）で進捗を一覧できる
\end{itemize}
\textbf{想定読者}：別冊『チームで安全に開発する』を終えた人。ブランチとPRは分かっている前提で書く。
\end{draftbox}

\begin{hintbox}
\textbf{この章の読み方}
\begin{itemize}[leftmargin=1.4em, nosep]
  \item まず動かしたい人 → 「Issueを立てる」→「IssueとPRをつなぐ」だけ読めばOK
  \item チームの運用ルールを決めたい人 → 「運用の型」へ
\end{itemize}
\end{hintbox}

% ---------------------------------------------------------------------
\section{Issueとは何か}

\begin{draftbox}
\textbf{書くこと}
\begin{itemize}[leftmargin=1.4em, nosep]
  \item Issue＝\textbf{リポジトリに貼る「作業の付箋」}。バグ報告・機能要望・やることメモを1件1枚で管理する
  \item PRとの違い：\textbf{Issueは「やること」、PRは「やった変更」}。Issueが先、PRが後
  \item チャットやメールで作業を頼むのとの違い：\textbf{リポジトリに残り、検索でき、変更履歴とひも付く}
  \item シリーズ共通の倉庫のたとえに接続する：Issueは「\textbf{倉庫に貼る作業指示書}」
\end{itemize}
\textbf{入れる図}：Issue（やること）→ ブランチ → PR（やった変更）→ マージ、の流れ図。
TikZ の \cmd{flownode} / \cmd{flowarrow} を使い、既存の図と体裁を揃える。
\end{draftbox}

% ---------------------------------------------------------------------
\section{Issueを立てる}

\begin{draftbox}
\textbf{書くこと}：GitHubのIssuesタブ →「New issue」→ タイトルと本文を書いて作成、までの手順。
\textbf{必要なスクリーンショット}（\cmd{images/04\_advanced/} に置く）
\begin{itemize}[leftmargin=1.4em, nosep]
  \item \cmd{issue\_new\_button.png}：Issuesタブと「New issue」ボタン
  \item \cmd{issue\_form.png}：タイトル・本文の入力画面
  \item \cmd{issue\_created.png}：作成後のIssue画面（番号 \cmd{\#12} が付くことを示す）
\end{itemize}
\end{draftbox}

\subsection{タイトルと本文の書き方}

\begin{draftbox}
\textbf{書くこと}
\begin{itemize}[leftmargin=1.4em, nosep]
  \item タイトルは\textbf{「何をどうしたいか」が一言で分かる形}にする（例：\cmd{ログイン画面のエラー表示が出ない}）
  \item 本文に書く3点：\textbf{現状／期待する状態／再現手順（またはやること）}
  \item 悪い例と良い例を並べて示す（「バグ直す」→ 何のバグか分からない）
  \item コミットメッセージの書き方（本編）と同じ考え方であることに触れる
\end{itemize}
\end{draftbox}

\subsection{Issueテンプレートを用意する}

\begin{draftbox}
\textbf{書くこと}：\cmd{.github/ISSUE\_TEMPLATE/} にテンプレートを置くと、New issue のときに雛形が出る。
付録A（別冊のPRテンプレート）と同じ形式で、\textbf{そのままコピーして使えるテンプレート}を載せる。

最低限、次の2種類を用意する：
\begin{itemize}[leftmargin=1.4em, nosep]
  \item \cmd{bug\_report.md}：不具合の報告用（現象／再現手順／期待する動作／環境）
  \item \cmd{task.md}：やること用（背景／完了条件／関連Issue）
\end{itemize}
\end{draftbox}

\begin{codebox}[markdown]
\begin{code}
（ここに bug_report.md の中身を貼る）
---
name: バグ報告
about: 動かない・想定と違う挙動を報告する
labels: bug
---

## 現象

## 再現手順

## 期待する動作

## 環境
\end{code}
\end{codebox}

% ---------------------------------------------------------------------
\section{ラベル・担当者・マイルストーンで整理する}

\begin{draftbox}
\textbf{書くこと}：Issueが増えたときに迷子にならないための3つの道具。

\begin{itemize}[leftmargin=1.4em, nosep]
  \item \textbf{ラベル}：種類で色分けする（\cmd{bug} / \cmd{enhancement} / \cmd{question} / \cmd{good first issue}）。おすすめの初期セットを表で示す
  \item \textbf{担当者（Assignees）}：誰がやるかを1人に決める。空のままにしない
  \item \textbf{マイルストーン}：締め切りやリリース単位でまとめる（例：\cmd{v1.0 提出}）
\end{itemize}
\textbf{入れる表}：ラベル名・色・使う場面の3列。
\textbf{必要なスクリーンショット}：\cmd{issue\_labels.png}（サイドバーでラベル・担当者を設定している画面）。
\end{draftbox}

% ---------------------------------------------------------------------
\section{IssueとPRをつなぐ}

\begin{draftbox}
\textbf{この節がこの章の山場}。PR本文に \cmd{Closes \#12} と書くと、\textbf{マージと同時にIssueが自動で閉じる}ことを示す。

\textbf{書くこと}
\begin{itemize}[leftmargin=1.4em, nosep]
  \item 使えるキーワード（\cmd{Closes} / \cmd{Fixes} / \cmd{Resolves}）と、単に参照だけしたいときの \cmd{\#12}
  \item ブランチ名にIssue番号を入れる運用（例：\cmd{feature/12-add-login}）
  \item Issue画面から直接ブランチを作る方法（\textbf{Create a branch} リンク）
  \item 「Issueを立てる → ブランチを切る → PRを出す → マージでIssueが閉じる」の一周を、実際の画面で通す
\end{itemize}
\textbf{入れる図}：この一周のフロー図（別冊 第1章のPRフロー図と体裁を揃える）。
\textbf{必要なスクリーンショット}：\cmd{issue\_linked\_pr.png}（PR本文の \cmd{Closes \#12} と、Issue側に出るリンク表示）。
\end{draftbox}

% ---------------------------------------------------------------------
\section{Projectsでボードにする}

\begin{draftbox}
\textbf{書くこと}：Issueを「未着手／作業中／レビュー中／完了」の列に並べて、進捗を一覧する。

\begin{itemize}[leftmargin=1.4em, nosep]
  \item Projects（ボード）の作り方と、Issueの追加方法
  \item ステータスの自動更新（PRがマージされたら「完了」に動かす設定）
  \item 少人数なら\textbf{ラベル運用だけでも足りる}こと。Projectsは人数と件数が増えてから入れるのがおすすめ、と\textbf{使いどころを限定}して書く
\end{itemize}
\textbf{必要なスクリーンショット}：\cmd{project\_board.png}（ボード全体）。
\end{draftbox}

% ---------------------------------------------------------------------
\section{運用の型}

\begin{draftbox}
\textbf{書くこと}：そのまま真似できる運用ルールを、別冊 付録Aと同じ調子で示す。

\begin{itemize}[leftmargin=1.4em, nosep]
  \item \textbf{作業を始める前に必ずIssueを立てる}（口頭で決めたことも1行で残す）
  \item \textbf{1 Issue = 1 ブランチ = 1 PR} を基本にする
  \item 週に1度、\textbf{開いたままのIssueを棚卸し}する（閉じ忘れの防止）
  \item 議論はチャットではなく\textbf{Issueのコメントに書く}（後から経緯を追えるようにする）
\end{itemize}
\end{draftbox}

% ---------------------------------------------------------------------
\section*{この章のまとめ}
\addcontentsline{toc}{section}{この章のまとめ}

\begin{draftbox}
本文を書き終えてから、次の骨子でまとめを作る。
\begin{itemize}[leftmargin=1.4em, nosep]
  \item Issueは「やること」、PRは「やった変更」
  \item Issueは\textbf{タイトルで内容が分かる}ように書く
  \item \cmd{Closes \#12} でPRとIssueをつなぐと、マージで自動的に閉じる
  \item 1 Issue = 1 ブランチ = 1 PR
\end{itemize}
次章は、この流れに\textbf{自動チェック}を足す話（GitHub Actions）につなげる。
\end{draftbox}
